\documentclass[12pt]{article}
\usepackage[T1]{fontenc}
\usepackage[utf8]{inputenc}
\usepackage{lmodern}
\usepackage{textcomp}
\usepackage{enumitem}
\usepackage{geometry}
\usepackage{graphicx}
\usepackage{amsmath, amssymb}
\geometry{margin=1in}
\begin{document}
\begin{center}
Chemistry - Undergraduate Level Assessment 15 \
Total Marks: 40 \
Answer all questions.
\end{center}

\section*{Multiple Choice (10 marks)}
\begin{enumerate}[label=\arabic*.]
\item The solid-state structures of the principal allotropes of elemental boron are made up of which of the following structural units?
\begin{enumerate}[label=(\alph*)]
    \item B$_{12}$ icosahedra
    \item B$_{8}$ cubes
    \item B$_{6}$ octahedra
    \item B$_{4}$ tetrahedra
\end{enumerate}
\item What is the orbital angular momentum quantum number, l, of the electron that is most easily removed when ground-state aluminum is ionized?
\begin{enumerate}[label=(\alph*)]
    \item 3
    \item 2
    \item 1
    \item 0
\end{enumerate}
\item Infrared (IR) spectroscopy is useful for determining certain aspects of the structure of organic molecules because
\begin{enumerate}[label=(\alph*)]
    \item all molecular bonds absorb IR radiation
    \item IR peak intensities are related to molecular mass
    \item most organic functional groups absorb in a characteristic region of the IR spectrum
    \item each element absorbs at a characteristic wavelength
\end{enumerate}
\item Reduction of D-xylose with NaBH 4 yields a product that is a
\begin{enumerate}[label=(\alph*)]
    \item racemic mixture
    \item single pure enantiomer
    \item mixture of two diastereomers in equal amounts
    \item meso compound
\end{enumerate}
\item Which of the following is a primary standard for use in standardizing bases?
\begin{enumerate}[label=(\alph*)]
    \item Ammonium hydroxide
    \item Sulfuric acid
    \item Acetic acid
    \item Potassium hydrogen phthalate
\end{enumerate}
\end{enumerate}

\section*{True / False (10 marks)}
\textit{Answer True or False}
\begin{enumerate}[label=\arabic*.]
\setcounter{enumi}{5}
\item True or False: At standard temperature and pressure, Br$_{2}$ is a red-brown volatile liquid.
\item True or False: BeCl 2 is more likely to act as a Lewis acid than SCl 2 because it has a higher electronegativity.
\item True or False: Protons and neutrons possess orbital angular momentum only.
\item True or False: Tetramethylsilane is used as a reference because its protons are not strongly shielded in NMR spectroscopy.
\item True or False: The normal modes of a carbon dioxide molecule that are infrared-active include symmetric stretching only.
\end{enumerate}

\section*{Long Form Response (20 marks)}
\textit{Answer each question with a clear and complete explanation.}
\begin{enumerate}[label=\arabic*.]
\setcounter{enumi}{10}
\item Discuss the factors that contribute to the limiting high-temperature molar heat capacity at constant volume (C\_V) of a gas-phase diatomic molecule, including its derivation and implications.
\item Discuss the implications of the Heisenberg uncertainty principle for a quantum mechanical particle in the lowest energy level of a one-dimensional box, focusing on position and momentum.
\end{enumerate}

\vfill
\noindent\textit{End of Paper}
\end{document}